\documentclass[11pt]{article}
\usepackage[margin=25mm]{geometry}
\usepackage{amsmath,amssymb}
\usepackage[hidelinks]{hyperref}
\title{Technical relations for establishing sensor contributions in multimodal AI}
\author{}
\date{29 October 2026}
\begin{document}
\maketitle
\noindent Editable equation companion to the v18 manuscript revision. The
repository file \texttt{docs/derivations.md} supplies the expanded explanations,
implementation mapping and evidence boundaries. These relations are analytical;
they are separate from attributed article values and participant-level evidence.

\section{Clock registration and sensor responses}
\[
t_{\mathrm{ref}}=a\,t_{\mathrm{device}}+b,\qquad
\min_{a,b}\sum_j\lambda_j[v_j-(a u_j+b)]^2.
\]
Times $u_j,v_j$ and $b$ are in seconds, $a$ is dimensionless, and
$\lambda_j\geq0$ is an event weight. This least-squares fit treats device event
times as sufficiently precise explanatory coordinates; appreciable errors in
both coordinates call for an errors-in-variables model. Identifying a rate and
offset requires at least two distinct event times with positive weights. Retained physical
events assess residual registration error after fitting. The response operators
labelled $H_x,H_z$ in the figure describe sensor physics, separately from the
timing-retention scalar $H$ below. The EEG--fNIRS conceptual relations are
\[
e(t)=L j(t)+\varepsilon_e(t),\qquad
c(t)=K[h*n](t)+u_s(t)+\varepsilon_c(t).
\]
Here $e$ is scalp voltage, $j$ source currents, $L$ the electrical lead field,
and $\varepsilon_e$ residual contributions. The approximate optical relation
uses relative haemoglobin estimates $c$, conceptual neural drive $n$,
haemodynamic kernel $h$, convolution $*$, effective sensitivity $K$,
superficial/systemic contributions $u_s$, and residual $\varepsilon_c$.
Units follow the declared measurement conversion and scaling. Clock correction
preserves the physiological response delay. These forward models are explanatory,
rather than implemented fitting procedures in the package.

\section{Relative timing error}
For residual errors expressed on a common timebase,
\[
\varepsilon_\Delta=\varepsilon_z-\varepsilon_x,\qquad
\sigma_\Delta^2=\sigma_z^2+\sigma_x^2
-2\mathrm{Cov}(\varepsilon_x,\varepsilon_z).
\]
Equal independent device standard deviations $s$ give
$\sigma_\Delta=\sqrt{2}s$. Identical common-mode errors cancel.

\section{Gaussian characteristic-function derivation}
Let $f$ be in hertz, and let the relative shift, in seconds, be independent of
the underlying feature:
\[
\varepsilon_\Delta\sim\mathcal N(\mu_\Delta,\sigma_\Delta^2),\qquad
Z=e^{i2\pi f\varepsilon_\Delta}.
\]
For a standard normal variable $U$, define
\[
F(v)=\mathbb E[e^{ivU}]
=\frac{1}{\sqrt{2\pi}}\int_{-\infty}^{\infty}e^{ivu}e^{-u^2/2}\,du.
\]
Integration by parts gives $F'(v)=-vF(v)$ and $F(0)=1$, hence
$F(v)=e^{-v^2/2}$. Therefore
\[
\mathbb E[Z]=e^{i2\pi f\mu_\Delta}
e^{-(2\pi f\sigma_\Delta)^2/2},\qquad
H(f,\sigma_\Delta)=|\mathbb E[Z]|
=e^{-(2\pi f\sigma_\Delta)^2/2}.
\]
$i$ is the imaginary unit; $\mu_\Delta$ is the mean relative shift. The mean
rotates phase, while variation reduces ensemble coherent amplitude. The
illustration uses $\mu_\Delta=0$. The distinct endpoints are
\[
|\mathbb E Z|=H,\qquad |\mathbb E Z|^2=H^2,\qquad
\mathbb E|Z|^2=1.
\]
The public helpers accept milliseconds and divide by 1000 explicitly.

\section{Finite-sample averages}
For $N$ independent, identically distributed unit phasors of any marginal
distribution, write $H=|\mathbb E Z|$. Then
\[
\bar Z_N=N^{-1}\sum_{j=1}^N Z_j,\qquad
\mathbb E|\bar Z_N|^2
=\frac{N+N(N-1)H^2}{N^2}
=H^2+\frac{1-H^2}{N}.
\]
The diagonal terms equal one; independence factorises each off-diagonal term
into $H^2$. This second moment differs from the squared ensemble mean and
from the mean sample magnitude. Correlated shifts instead require
\[
\mathbb E|\bar Z_N|^2
=N^{-2}\sum_{j,k}\mathbb E\!\left[
e^{i2\pi f(\varepsilon_{\Delta,j}-\varepsilon_{\Delta,k})}\right].
\]
Gaussian marginal distributions determine $H$; sampling dependence determines
convergence and uncertainty of empirical averages.

\section{A chosen timing budget}
For $f>0$ and $0<\eta\leq1$, solving $H\geq\eta$ gives
\[
\sigma_\Delta\leq\frac{\sqrt{-2\ln\eta}}{2\pi f}.
\]
At the illustrative amplitude-retention criterion $\eta=0.9$, this is
3.65295294 ms at 20 Hz and 1.82647647 ms at 40 Hz. The corresponding coherent
power fraction is $0.9^2=0.81$. A 90\% coherent-power requirement uses
$H\geq\sqrt{0.9}$. These are chosen analytical criteria, rather than universal
instrument thresholds. Carrier phase, envelopes, finite-window features,
haemodynamic dynamics and task performance require their own evaluations.

\section{Affine expansion in fixed features}
\[
x\in\mathbb R^{d_x},\quad z\in\mathbb R^{d_z},\quad
s=w_x^\top x+w_z^\top z+b
=[w_x;w_z]^\top[x;z]+b.
\]
For one unconstrained scalar output with fixed features, the total count is
$d_x+d_z+1$, compared with $d_x+1$ for the existing-input score. The added block
introduces $d_z$ coefficients. Joint fitting can also change $w_x$ and $b$.
An affine score followed by a logistic link
$p=(1+e^{-s})^{-1}$ yields a probability nonlinear in the declared features.
Dimension bookkeeping alone establishes neither target information nor benefit.

\section{Interactions, encoders and fusion}
\[
r(x,z)=x^\top A z,\quad A\in\mathbb R^{d_x\times d_z},\qquad
r(c x,c z)=c^2r(x,z).
\]
The interaction is linear in either input with the other held fixed, generally
nonlinear in $[x;z]$, and linear in the $d_xd_z$ entries of $A$ after declaring
product features. An unconstrained bilinear block plus the affine $z$ block
adds $d_xd_z+d_z$ coefficients. Architectural constraints change that count.
More generally,
\[
h_x=E_{\phi_x}(x),\quad h_z=E_{\phi_z}(z),\quad
s=g_\theta(h_x,h_z,q,m).
\]
$\phi_x,\phi_z$ are encoder parameter blocks; $h_x,h_z$ have dimensions
$k_x,k_z$; $\theta$ parameterises fusion; $q$ records measured quality; and $m$
records stream availability. A declared concatenating readout has input
dimension $k_x+k_z+\dim(q)+\dim(m)$. Architecture and update schedule determine
trainable parameter counts. An affine head is affine in its fixed
representations, while nonlinear encoders make the complete mapping nonlinear.

\section{Interpretation}
Supervision describes the origin of training targets, and one system can
combine supervised, self-supervised and other unsupervised stages. The added
observation's role determines its evidence requirement: an audit distinguishes
a named explanation, a correction receives interference and target-preservation
checks, and a prediction input receives a matched comparison at a common
decision deadline. These analytical identities alone establish neither causal
identification nor physiological, predictive or clinical benefit.
\end{document}
